% !TeX spellcheck = ru_RU
% !TEX root = vkr.tex

\section{Обзор существующих решений}
\label{sec:relatedworks}

Настоящая работа посвящена реорганизации инфраструктуры для измерения производительности графовых алгоритмов на основе взаимодействующих сетей~\cite{inpla}. Под инфраструктурой здесь понимается совокупность скриптов, настроек и соглашений, автоматизирующих подготовку тестовых матриц, запуск измерений и сбор результатов для набора инструментов.

В качестве объектов сравнения рассматриваются реализации на основе взаимодействующих сетей: \emph{QTreeInpla}~\cite{inpla}, \emph{QTreeVine}~\cite{vine}, а также однопоточная реализация на \emph{F\#}, являющаяся их прототипом. Кроме того, используются эталонные инструменты: \emph{NetworkX}~\cite{networkx} и \emph{LAGraph}~\cite{lagraph}, построенный поверх \emph{SuiteSparse:GraphBLAS}~\cite{graphblas}.

Измерения проводятся для трёх алгоритмов: обхода в ширину (\emph{BFS}), поиска кратчайших путей из источника (\emph{SSSP}) и подсчёта треугольников (\emph{Triangle Counting})~\cite{trianglecounting}.

Ранее инфраструктура состояла из набора отдельных скриптов в каталоге \texttt{scripts/}. Такая организация имела ряд недостатков: логика загрузки и подготовки матриц дублировалась между скриптами, набор матриц и инструментов не был описан в одном месте, добавление нового инструмента требовало правок в нескольких файлах, а результаты измерений не имели единого формата и могли оказаться несовместимы с кодом построения графиков.

Реорганизация была направлена на устранение этих недостатков: единый интерфейс запуска, общие модули, формализованные правила подключения инструментов, автоматизация работы с матрицами и единое соглашение о хранении результатов.
